\documentclass{IQ_Thesis}%Inicia el tipo de documento, tamaño de letra, tamaño de papel, doble cara, tipo artículo.


\title{Título del reporte}%Título del documento, no necesariamente se desplega.
\author{Nombre del autor}%Nombre del autor del documento, no necesariamente se desplega.
\date{Fecha}%Fecha de cuando se realiza el documento, no necesariamente se desplega.

% In your preamble 
 
\makeatletter
\newcommand*{\addFileDependency}[1]{% argument=file name and extension
  \typeout{(#1)}
  \@addtofilelist{#1}
  \IfFileExists{#1}{}{\typeout{No file #1.}}
}
\makeatother
 
\newcommand*{\myexternaldocument}[1]{%
    \externaldocument{#1}%
    \addFileDependency{#1.tex}%
    \addFileDependency{#1.aux}%
}
% In your preamble
 
\myexternaldocument{D._METHODOLOGY/METHOD.tex}
%-----------INICIO DEL DOCUMENTO-----------%
\begin{document}%Inicia el documento.
\thispagestyle{empty}%Elimina número de página.

%-----------PORTADA-----------%
\subfile{A._COVER/PAGE}


%-----------RESUMEN-----------%
\myresume{0}{RESUMEN}%Inicia resumen
\subfile{B._RESUME/ABSTRACT}


%-----------ÍNDICE-----------%
%Los cuadros rojos que se observan en el pdf son hipervínculos que se añaden al fdp al compilarlo.
\blankpage
\begin{center}
\centering{\tableofcontents}%Tabla de contenidos.
\end{center}

%-----------ÍNDICE DE CUADROS-----------%
\blankpage
\begin{center}
\listoftables
\end{center}
\newpage
\thispagestyle{empty}
%-----------ÍNDICE DE FIGURAS-----------%
\begin{center}
\listoffigures
\end{center}






%----------MARCO TEÓRICO----------%

\myfundaments{1}{CAPÍTULO 1 \newline INTRODUCCIÓN}
\pagenumbering{arabic}%Numeración arábiga.
\setcounter{page}{1}%Iniciar en 1.

\subfile{C._FUNDAMENTS/CAP_I}
\subfile{C._FUNDAMENTS/CAP_II}
\subfile{C._FUNDAMENTS/CAP_III}
\subfile{C._FUNDAMENTS/CAP_IV}
%\subfile{C._FUNDAMENTS/CAP_V}



%----------METODOLOGÍA----------%
\mymethod{5}{METODOLOGÍA EXPERIMENTAL}
\subfile{D._METHODOLOGY/METHOD}
  		
%----------ANÁLISIS DE RESULTADOS----------%
\myresults{6}{ANÁLISIS DE RESULTADOS}
\subfile{E._RESULTS/ANALYSIS}




%----------CONCLUSIONES Y RECOMENDACIONES----------%
\myconclusions{7}{CONCLUSIONES Y RECOMENDACIONES}
\subfile{F._CONCLUSIONS/RECOMENDATIONS}




%----------NOMENCLATURA----------%
\subfile{G._NOMENCLATURE/SYMBOLS}
\pagestyle{fancy}
\printnomenclature[20 mm]
\pagestyle{fancy}


%----------BLIBLIOGRAFÍA----------%
\pagestyle{fancy}
\bibliographystyle{apalike-es}%Tipo de referenciado.
\bibliography{Referencias.bib}%Invoca al documento de ese nombre.
%El ambiente no permite desplegar referencias sin haber sido citadas con antelación.
\pagestyle{fancy}










%----------APÉNDICES----------%
\myappendix{9}{APÉNDICES}
\pagestyle{fancy}
\subfile{H._DATA/EXPERIMENTAL}

%----------ANEXOS----------%
\myannexed{10}{ANEXOS}
\subfile{I._APPENDIX/ANNEXED}
\clearpage
%----------GOALS CHECK LIST----------%
%\subfile{J._CHECK_LIST/GOALS}

\end{document}
%----------ANTIGUA NOMENCLATURA----------%
%\begin{table}[!hbt]
%		\begin{center}
%		\begin{tabularx}{140mm}{X l l}%Tabularx permite adaptar las columnas con X al ancho especificado haciéndolas más grandes o más pequeñas.
        
%----------Mayúsculas----------%
%$A$&Área normal a la dirección del fluido &m$^{2}$ \\
%$E$&Energía & J\\

%&&\\ %Fila vacía en la Tabla
%----------Minúsculas----------%
%$c$&Diferencia longitudinal de alambre fusible& m\\
%$e$&Correción energética por alambre fusible & J\\

%&&\\
%----------Griegas----------%
%$\alpha$&Correción de energía cinética& adim\\
%$\epsilon$&Emisividad& adim\\
%&&\\ %Fila vacía en la Tabla
%\textbf{Subíndices}&\\

%$\infty$&Indica muy lejano\\
%$+$&Indica adición de reactivo o mezcla\\

%&\\
%\textbf{Superíndices}&\\
%$\cdot$&Indica por unidad de tiempo\\
%$\rightarrow$&Indica cantidad vectorial\\

%	\end{tabularx}
%		\end{center}
%	\end{table}







